\documentclass[../main.tex]{subfiles}
\begin{document}

\section{Problem Statement}
\label{sec:dvd}
Large language models achieve strong results across natural-language tasks, yet their
internal computation remains largely opaque. Mechanistic interpretability seeks to
reverse-engineer that computation into human-understandable algorithms rather than only
relating inputs to outputs. A central goal of this effort is circuit tracing, which
identifies the pathways of activation inside a model that implement a particular behaviour.
Recent work makes this concrete: Ameisen et al.~\cite{ameisen2025circuit} use a cross-layer
transcoder as a replacement model to extract an \emph{attribution graph}, a directed graph
whose nodes are interpretable features and whose edges are the linear effects between them,
recording the steps the model takes to produce the next token for a given prompt. The
attribution graph is a faithful, prompt-specific account of the model's computation, and
it is the object this thesis studies.

The attribution graph, however, is too large to read. A single prompt induces thousands of
feature nodes joined by a dense web of edges, and making sense of one graph can take an
expert hours, or becomes infeasible for long prompts. The graph must therefore be reduced
to a compact form before a human can use it, but the reduction must not distort the
computation it represents. This thesis addresses the resulting problem: how to summarize an
attribution graph automatically into a small, human-readable graph that remains faithful to
the model's behaviour.

\section{Background and Problems of Research}
\label{sec:giaiphap}
The circuit-tracing methodology already provides a way to shrink an attribution graph, but
it relies on human judgment at every step. After building the graph, Ameisen et
al.~\cite{ameisen2025circuit} prune it by keeping the nodes and edges that account
for most of the influence on the output and discarding the rest, and then group the
surviving features into \emph{supernodes} by inpsection, collecting features that play the
same role into a single labelled unit so the circuit can be drawn as a small diagram. Both
steps are manual: pruning thresholds are tuned by hand, and grouping a single circuit can
take hours of expert effort, so the analysis does not scale beyond a handful of prompts.

Related efforts address other parts of the problem but leave the structural reduction open.
Work on circuits in the neuron
basis~\cite{arora2026sparse} avoids transcoder training and reconstruction error, yet the
authors note that the resulting circuits are still too large to read and call for principled
grouping methods.
ADAG~\cite{arora2026adag} addresses this by introducing attribution profiles to quantify the role of 
neurons in the graph and use spectral clustering to cluster them. This approach does not take account for structral information of the neurons,
making the clustering less mechanisticly faithful. No existing method produces a compact and faithful structural summary
of a signed attribution graph by jointly pruning and clustering its feature nodes under an
explicit faithfulness objective.

\section{Research Objectives and Conceptual Framework}
This thesis aims to summarize attribution graphs automatically so that researchers can
understand a model's behaviour without manual reduction. The objective is a method that
takes an attribution graph and returns a small \emph{summary graph} of supernodes that is
interpretable, faithful to the original computation, and produced without human
intervention. The proposed solution responds directly to the limitations of the previous
section. In place of hand-tuned pruning, it scores every node and edge by a combination of
its downstream influence on the output and its upstream relevance to the input, and keeps
only the important ones. In place of grouping features by inspection, it clusters them by a
role similarity under an explicit objective that rewards atomic supernodes while preserving
the causal edges between them. Finally, it aggregates the surviving edges into an acyclic,
labelled supernode graph, so the output is a directed acyclic graph that a reader can follow
at a glance. Each stage replaces one manual step with an automated, objective-driven
procedure.

\section{Contributions}
This thesis makes three contributions:
\begin{itemize}[leftmargin=*,topsep=2pt,itemsep=2pt]
\item It formalizes the problem of summarizing an attribution graph, defining the summary
graph together with explicit desiderata --- importance retention, atomicity,
causal-information preservation, and interpretability --- and an objective that makes them
measurable.
\item It proposes an automated three-stage pipeline --- pruning, clustering, and
interpreting --- that produces a concise summary graph that is faithful to the original attribution graph.
\item It evaluates the method on the BATS analogy and multi-hop reasoning datasets,
showing through steering interventions that the resulting summaries are causally faithful
and that the supernodes behave as the single mechanisms.
\end{itemize}

\section{Organization of Thesis}
The remainder of this thesis is organized as follows. Chapter~2 reviews the background the
method depends on, tracing the line from superposition and sparse dictionary learning to the
cross-layer transcoders and circuit-tracing pipeline that produce attribution graphs, and
surveys the related work on describing and summarizing circuits from which the thesis takes
its motivation. Chapter~3 presents the proposed method, formalizing the summary-graph
problem and detailing its three stages of pruning, clustering, and interpreting. Chapter~4
analyzes the method's theoretical properties, in particular how reconstruction error limits
the influence score on which pruning relies. Chapter~5 evaluates the method empirically,
defining the faithfulness metrics and reporting the steering-intervention study that
validates the summary graphs against the underlying model. Chapter~6 discusses the findings
and their limitations and concludes with directions for future work.

\end{document}
